\documentclass[11pt,a4paper]{article}
\usepackage[utf8]{inputenc}
\usepackage[vietnamese]{babel}
\usepackage{amsmath,amssymb}
\usepackage{geometry}
\geometry{a4paper, total={180mm,260mm}, left=15mm, right=15mm, top=18mm, bottom=18mm}

\usepackage{tikz}
\usetikzlibrary{shapes.geometric, arrows.meta, positioning, calc}
\tikzset{
  arr/.style={-{Stealth[scale=0.8]}},
  every node/.append style={align=center}
}
\usepackage{tcolorbox}
\tcbuselibrary{skins,breakable}
\usepackage{fontawesome}
\usepackage{fancyhdr}
\usepackage{enumitem}

\pagestyle{fancy}
\fancyhf{}
\renewcommand{\headrulewidth}{0.5pt}
\renewcommand{\footrulewidth}{0.5pt}
\fancyhead[L]{\small\bfseries\color{violet!85!black}\faBook\ SỔ TAY KIẾN THỨC TOÁN 2}
\fancyhead[R]{\small\bfseries\color{violet!85!black}Lớp Toán Thầy Lĩnh -- loptoanthaylinh.vn}
\fancyfoot[C]{\small\color{gray}-- Sổ tay kiến thức: Bài 17. Thực hành cân đo ki-lô-gam và lít --}

\setlength{\parindent}{0pt}
\setlength{\parskip}{0pt}

\begin{document}

% --- BANNER BÀI HỌC ---
\begin{tcolorbox}[
  enhanced,
  colback=violet!8!white,
  colframe=violet!70!black,
  arc=3mm, boxrule=1.2pt,
  left=3mm, right=3mm, top=2mm, bottom=2mm,
  halign=center
]
  {\LARGE\bfseries\color{violet!90!black}\faBookmark\ BÀI 17: THỰC HÀNH VÀ TRẢI NGHIỆM VỚI KI-LÔ-GAM, LÍT}\\[1mm]
  {\itshape\footnotesize\color{black!75} Kĩ năng thực tế: Đọc vạch cân đồng hồ, quan sát mức nước ca đong và ước lượng trên 1 trang A4}
\end{tcolorbox}

\vspace{-2mm}

\noindent
\begin{minipage}[t]{0.49\textwidth}\vspace{0pt}
\begin{tcolorbox}[
  enhanced,
  colback=violet!4!white,
  colframe=violet!75!black,
  arc=2mm, boxrule=0.8pt,
  left=2mm, right=2mm, top=1.8mm, bottom=1.8mm,
  title={\footnotesize\bfseries\color{white}\faKey\ 1. Kĩ năng đọc cân đồng hồ}
]
{\footnotesize
\begin{itemize}\setlength{\itemsep}{1pt}\setlength{\leftskip}{-2.5mm}
  \item \textbf{Kim chỉ vị trí nào:} Đọc số tại vạch mà kim đồng hồ chỉ vào.
  \item Khi đặt vật lên đĩa cân, kim quay theo chiều kim đồng hồ từ mốc số $0$.
  \item \textbf{Ví dụ:} Kim chỉ vào số $3$ $\implies$ Vật nặng $3\text{ kg}$.
  \item Kim chỉ ở giữa số $2$ và số $3$ $\implies$ Nặng hơn $2\text{ kg}$ và nhẹ hơn $3\text{ kg}$.
\end{itemize}
}
\end{tcolorbox}
\end{minipage}
\hfill
\begin{minipage}[t]{0.49\textwidth}\vspace{0pt}
\begin{tcolorbox}[
  enhanced,
  colback=teal!4!white,
  colframe=teal!75!black,
  arc=2mm, boxrule=0.8pt,
  left=2mm, right=2mm, top=1.8mm, bottom=1.8mm,
  title={\footnotesize\bfseries\color{white}\faSliders\ 2. Kĩ năng đong nước và ước lượng}
]
{\footnotesize
\begin{itemize}\setlength{\itemsep}{1pt}\setlength{\leftskip}{-2.5mm}
  \item Đặt ca đong trên mặt bàn phẳng để quan sát chính xác mặt thoáng của nước.
  \item Mặt nước dâng đến vạch chia nào thì lượng nước đúng bằng số ghi ở vạch đó ($1\ l, 2\ l,\dots$).
  \item \textbf{Ước lượng đời sống:}\\
  - Cặp sách học sinh: Khoảng $2\text{ kg} - 3\text{ kg}$.\\
  - Con gà con vịt: Khoảng $1\text{ kg} - 2\text{ kg}$.\\
  - Chai dầu ăn gia đình: Thường là $1\ l$ hoặc $2\ l$.
\end{itemize}
}
\end{tcolorbox}
\end{minipage}

\vspace{-1.5mm}
\begin{tcolorbox}[
  enhanced,
  colback=orange!4!white,
  colframe=orange!80!black,
  arc=2mm, boxrule=0.8pt,
  left=2mm, right=2mm, top=1.5mm, bottom=1.5mm,
  title={\footnotesize\bfseries\color{white}\faSitemap\ 3. Bảng phân loại ứng dụng: Cân nặng vs Dung tích}
]
\centering
\begin{tikzpicture}[x=1.5cm, y=0.8cm]
  \node[rectangle, rounded corners=2pt, draw=violet!80!black, fill=violet!15, text width=4.5cm, align=center, font=\scriptsize\bfseries] at (0,0) {
    \textbf{ĐO KHỐI LƯỢNG (kg)}\\Dụng cụ: Cân đĩa, cân đồng hồ\\Vật mẫu: Gạo, bí ngô, hoa quả, thịt
  };
  \node[rectangle, rounded corners=2pt, draw=teal!80!black, fill=teal!15, text width=4.5cm, align=center, font=\scriptsize\bfseries] at (4.5,0) {
    \textbf{ĐO DUNG TÍCH (l)}\\Dụng cụ: Ca đong, bình định mức\\Chất lỏng: Nước lọc, sữa tươi, mật ong, dầu
  };
\end{tikzpicture}
\end{tcolorbox}

\vspace{-1.5mm}
\noindent
\begin{minipage}[t]{0.49\textwidth}\vspace{0pt}
\begin{tcolorbox}[
  enhanced,
  colback=green!4!white,
  colframe=green!60!black,
  arc=2mm, boxrule=0.8pt,
  left=2mm, right=2mm, top=1.8mm, bottom=1.8mm,
  title={\footnotesize\bfseries\color{white}\faCheckCircle\ 4. Ví dụ thực hành chuẩn mực}
]
{\footnotesize
\textbf{1. Đọc số đo thực tế:}\\
- Cân một quả dưa hấu, kim chỉ số $4$ $\implies$ Quả dưa hấu nặng $\mathbf{4\text{ kg}}$.\\
- Rót đầy $3$ ca nước loại $2\ l$ vào xô $\implies$ Trong xô có: $2 + 2 + 2 = \mathbf{6\ l}$ nước.
\textbf{2. Điền đơn vị thích hợp (kg hay l):}\\
- Con chó nặng khoảng $8$ \textbf{kg}.\\
- Bình xăng xe máy chứa khoảng $4$ \textbf{l}.
}
\end{tcolorbox}
\end{minipage}
\hfill
\begin{minipage}[t]{0.49\textwidth}\vspace{0pt}
\begin{tcolorbox}[
  enhanced,
  colback=red!4!white,
  colframe=red!75!black,
  arc=2mm, boxrule=0.8pt,
  left=2mm, right=2mm, top=1.8mm, bottom=1.8mm,
  title={\footnotesize\bfseries\color{white}\faExclamationTriangle\ 5. Mẹo tránh bẫy trắc nghiệm}
]
{\footnotesize
\begin{itemize}\setlength{\itemsep}{1pt}\setlength{\leftskip}{-2.5mm}
  \item \textbf{Bẫy nhìn lệch vạch số $0$:}\\
  Trước khi cân, phải kiểm tra kim cân đã về đúng vạch số $0$ chưa (chỉnh cân về chuẩn).
  \item \textbf{Bẫy điền nhầm đơn vị:}\\
  Chất rắn $\to$ điền \textbf{kg}. Chất lỏng $\to$ điền \textbf{l}.
  \item Tránh đoán mò khối lượng vượt quá thực tế (ví dụ: con mèo nặng $50\text{ kg}$ là sai phi lý!).
\end{itemize}
}
\end{tcolorbox}
\end{minipage}


\end{document}
